% New admissibility argument for initializer-driven computation dispatch.
% Inclusion fragment; requires amsmath and amssymb.

\section{Using the initializer to select computation}
\label{sec:seed-dispatch}

The initialized scalar iterator has the usual equations
\[
 I(0,a;v.h)=a,\qquad
 I(Sn,a;v.h)=h[I(n,a;v.h)/v].
 \tag{I0, IS}
\]
The following restriction removes a computation opcode from the
certificate interface without changing these operations or the native
consequence relation.  It concerns the \emph{syntax} of an initializer,
not whether that initializer denotes a positive natural number.

\paragraph{A fixed nonempty representation of zero.}
Let
\[
               C:=I(S(S0),S0;v.0).
 \tag{Auxiliary zero}\label{eq:dispatch-zero-term}
\]
One successor computation, at predecessor $S0$, proves
\[
                              C=0.\label{eq:dispatch-zero}
 \tag{Zero certificate}
\]
The predecessor and initializer in this computation are both nonempty
canonical terms.  The update is the constant zero, so the right side of
the computation is exactly zero.  This construction does not use a
previous computation with initializer zero, a zero-predecessor repair,
an induction, or a presumed identity for a new numerical constructor.

\paragraph{Restricted computation system.}
In canonical numerical serialization, use scalar $0$ for zero and
nonempty arrays for successor and iteration.  Consider the same native
calculus with its direct computation leaves restricted as follows:

\begin{itemize}
\item I0 may be used directly only when its initializer is not literal
  zero.
\item IS may be used directly only when both its initializer and its
  supplied predecessor are not literal zero.
\end{itemize}

The predecessor restriction is already present in the inherited
predecessor-specified interface.  The new restriction is on the
initializer.  A parameter or a nonempty term is permitted even when its
value happens to be zero.  In particular, $C$ is permitted.  Keep the
ordinary native equality-transport rule and all other logical rules
unchanged.

\paragraph{Admissibility theorem.}
Every unrestricted I0 and IS instance has a finite derivation in this
restricted system.  Consequently the restricted and unrestricted
calculi prove exactly the same native sequents, with exactly the same
numerical and formula languages.

\paragraph{Proof.}
Instances satisfying the stated syntactic conditions are direct.
For a base computation with initializer zero, take the fresh-variable
equation motive
\[
                  E(w):=\bigl(I(0,w;v.h)=w\bigr).
\]
The restricted base equation proves $E(C)$.  Equality transport along
the independently proved $C=0$ gives $E(0)$, which is precisely the
required unrestricted base instance.

For a successor computation with permitted predecessor $n$ but zero
initializer, take
\[
 E(w):=\left(I(Sn,w;v.h)=h[I(n,w;v.h)/v]\right).
\]
The restricted successor equation proves $E(C)$, and the same transport
gives $E(0)$.  The displayed variable $w$ is fresh for $n$, $h$, and all
retained assumptions.  Only the indicated initializer occurrences are
substituted; unrelated zero occurrences or passive arguments in $h$
remain fixed.

For a zero predecessor but a permitted initializer $a$, use instead
\[
 E(w):=\left(I(Sw,a;v.h)=h[I(w,a;v.h)/v]\right).
\]
Again $E(C)$ is a permitted successor computation and transport gives
the desired instance $E(0)$.

If both predecessor and initializer are zero, a single shared motive
repairs them simultaneously:
\[
 E(w):=\left(I(Sw,w;v.h)=h[I(w,w;v.h)/v]\right).
\]
Both supplied values in $E(C)$ are nonempty, and transport along $C=0$
gives exactly $E(0)$.  Bound state variables are renamed before these
substitutions, so no capture occurs.

Each repaired computation leaf therefore uses one permitted computation
for $C=0$, one permitted computation for $E(C)$, and one transport
inference.  Replacing computation leaves in a native derivation proves
the inclusion from the unrestricted calculus to the restricted one;
the reverse inclusion is immediate.  Contexts and final sequents are
unchanged throughout.\hfill$\square$

This is a three-node repair when computation leaves and transport are
counted as one node each and the proof of $C=0$ is copied.  If a tree
derivation has $N$ nodes and $r$ excluded computation leaves, the stated
replacement has exactly $N+2r$ nodes.  If a DAG proof retains one shared
node proving the closed lemma $C=0$, the repaired part instead has
$2r+1$ nodes, giving $N+r+1$ distinct nodes when $r>0$.  A certificate
format with persistent proved assumptions can implement corresponding
sharing by supplying this closed computation once before the affected
continuation.  These are counts for the displayed construction, not
lower bounds on shortest certificates or claims about a particular
producer's serialization.

\paragraph{Why the opcode can disappear.}
Write a four-field packet as $[m,s,u,t]$.  The established elimination
branch uses a truthy major $m$.  In a computation packet, let $s$ be the
predecessor flag/value, $u$ the initializer, and $t$ the update.  A truthy
$u$ selects the computation branch.  The checker constructs the source
counter as $s$ when it is the zero flag and as $Ss$ otherwise, and
generates I0 or IS.  A paired-introduction packet puts zero in the unused
$u$ field.  Thus the dispatch can be written
\begin{verbatim}
    if m:
        ... elimination ...
    elif u:
        ... computation ...
    else:
        ... paired introduction ...
\end{verbatim}
The direct computation branch can no longer receive a literal-zero
initializer.  The admissibility theorem supplies exactly that missing
case.  No test that $u$ denotes a nonzero number is performed or needed.

\paragraph{The fresh-name obligation on noncanonical packets.}
The dispatch change admits computation packets with any false-valued
$m$, including the atomic-progress marker $m=0$.  In the common
continuation, that marker can cause the generated hypothesis to be
opened at the new rigid parameter $b$, rather than at literal zero.
It would be incorrect to assert that the resulting $H(b)$ is uniformly
true.  Use the existing zero-instance invariant instead:
\[
        \vdash H(\vec0),\qquad
        \vec b\cap\operatorname{FV}(\Gamma,A)=\varnothing.
\]
Here $\vec b$ includes the designated new integer name $b$ and every
additional negative fractional name in the newly admitted unit interval
that occurs in the finite raw certificate.  Such names must not be
silently identified with the designated integer.  Given a finite
derivation of $\Gamma,H(\vec b)\vdash A$, rename its later
eigenparameters if necessary, substitute zero for every member of
$\vec b$, and cut the resulting assumption $H(\vec0)$ against its
computation derivation.  This
gives $\Gamma\vdash A$.  The old-floor validation of the retained goal
and context is what establishes the displayed freshness condition.
Consequently the new raw packet possibilities are sound under the same
zero-completion and opening lemmas as the predecessor-specified kernel.
They do not justify silently strengthening a known zero instance into
a uniform computation theorem.

\paragraph{Interpretations and the scope of the improvement.}
The theorem preserves every native sequent through the identity map on
object expressions.  It therefore transfers the inherited whole-domain
bi-interpretation and equiconsistency with iPRA without a new domain,
new equality relation, or new primitive numerical operation.  The
contribution is a removable syntactic restriction and a data-driven
certificate dispatch.  It is an implementation simplification, not a
new arithmetic foundation or a proof of absolute minimality.

Relative to the inspected pinned draft collection, the new step is
using the initializer itself to discriminate computation packets and
proving that the resulting missing computation instances are admissible.
The broader technique of replacing an unavailable zero representation
by a proved nonempty representative is already used in the supplied
drafts, including their predecessor and witness repairs.  The present
argument adapts that attributed technique to eliminate the computation
opcode.  No priority claim outside the inspected corpus is made.
